\section{Complete fleets and the meaning of support}\label{sec:framework}
Fix mandatory services, directed movement opportunities, a charging horizon
and physical equipment. A complete plan $x\in\X$ covers every mandatory
service exactly once, assigns the services to used buses, routes those buses,
and schedules charging subject to battery and shared-resource constraints.
Uncounted buses neither exchange energy nor incur operating cost. When a
variant admits paid idle buses, they count as used participants under that
variant's fleet and access rules. Let $e(x)\in\R^T$ denote net grid energy
by market interval, positive for purchases and negative for exports, and let
$c(x)$ be intrinsic operating cost. Charge-only variants restrict
$e(x)\in\R_+^T$. Each used battery starts full and finishes at the
same required full inventory. Assume the projected image
\[
 S=\{(e(x),c(x)):x\in\X\}
\]
is nonempty and compact. All physical planning, price-response and hull
problems below use exactly this common set.

Let $F:\R^T\to\R\cup\{+\infty\}$ be proper, closed and convex, finite on
all feasible loads. It represents incremental supply cost conditional on
fixed background demand, on the declared feasible signed-load domain when
exports are allowed. Define
\begin{align}
 D&=\min_{x\in\X}\{c(x)+F(e(x))\},\label{eq:D}\\
 \CH&=\min_{(e,c)\in\operatorname{conv}S}\{c+F(e)\},
 \qquad\Delta=D-\CH\ge0.\label{eq:CH}
\end{align}
The hull mixes \emph{complete feasible plans}; it is not a duty-column
relaxation unless equivalence is separately proved. Its mean load need not
be executable as a physical daily schedule. Evaluating supply cost at that
mean also differs from averaging the costs of plans operated on different
days.

For a posted tariff $p\in\R^T$, the complete-fleet response and regret are
\begin{equation}\label{eq:regret}
 V(p)=\min_{y\in\X}\{c(y)+p^\top e(y)\},\qquad
 r(x;p)=c(x)+p^\top e(x)-V(p).
\end{equation}
The posted tariff is fixed across alternatives $y$. A strategic operator
anticipating a tariff reset would face a different objective and needs a
separate settlement model.

\begin{definition}[Weak own-load support]
A physical plan $x$ is supported if there exists an admissible marginal
price $p\in\partial F(e(x))$ for which $r(x;p)=0$. Alternative plans may
tie. If an economic convention restricts the price set, existence is tested
only in that declared subset.
\end{definition}
For differentiable supply we use its gradient. For our quadratic examples
this means the gradient of the unrestricted quadratic formula, including
at a zero load. In charge-only variants, orthant-domain normals are not additional
generation-price choices. With nonsmooth supply, a failed support test at
one chosen multiplier does not exclude support at another admissible price.

For the following classical equivalence, assume either
\[
 \operatorname{ri}\operatorname{conv}e(\X)
 \cap\operatorname{ri}\operatorname{dom}F\ne\varnothing,
\]
or that a finite convex minorant on $\R^T$ agrees with $F$ on all hull
loads. The latter applies to our quadratic. These assumptions ensure an
attained hull-dual price; a restricted economic price set must contain such
a price for the corresponding converse.

\begin{proposition}[Classical support criterion]\label{prop:support}
For every physical $x$ and $p\in\partial F(e(x))$,
\begin{equation}\label{eq:lower}
 r(x;p)\ge c(x)+F(e(x))-\CH\ge\Delta.
\end{equation}
Under the stated qualification, $D=\CH$ if and only if some physical plan
is supported. If $D=\CH$, every physical planner optimum is supported by
every \emph{admissible} attained maximizer of $V(p)-F^*(p)$, where
$F^*(p)=\sup_u\{p^\top u-F(u)\}$.
\end{proposition}
This specializes established nonconvex support and convexification
arguments \citep{oneill2005,gribikHoganPope2007,huaBaldick2017}; the proof
is included in Appendix~\ref{app:support} to make assumptions and quantifiers
explicit. It does not say that every subgradient supports a planner optimum,
or that a price-adjustment algorithm converges. Without a price restriction,
every attained dual maximizer has zero regret and is a subgradient at each
physical planner optimum. An external price restriction adds a separate
admissibility test to that mathematical conclusion.

\begin{remark}[Admissibility is a separate quantifier]
For a single physical point $(e,c)=(0,0)$ and $F(e)=|e|$, both values
are zero and every $p\in[-1,1]$ attains the hull dual. If the declared
admissible price set is $\{0\}$, the plan has admissible support at zero;
the dual maximizer $p=1$ remains mathematically supporting but is inadmissible.
Containing one dual maximizer does not mean containing them all.
\end{remark}

The participant in \eqref{eq:regret} controls all buses. Shared infrastructure
does not automatically define unilateral alternatives for separately owned
vehicles. Section~\ref{sec:ownership} introduces reserved rights before
making an individual-participant comparison.
